\chapter{Sheaves, Sheafification, and Modalities}\label{ch04:main}

This chapter completes the background half of the monograph. Chapter~2 fixed the
categorical vocabulary; Chapter~3 introduced sites: categories equipped with a
Grothendieck topology that declares which sieves count as covers. The present
chapter explains what covers are \emph{for}. A topology is interesting only
because it selects a class of presheaves---the sheaves---whose sections are
determined by compatible local data, and because the passage from an arbitrary
presheaf to its best sheaf approximation---sheafification---is a functor with
exceptionally good properties. Composing sheafification with the inclusion of
sheaves back into presheaves yields a left-exact idempotent monad, which this
monograph calls a \emph{modality}. Modalities are the objects whose existence,
distinctness, triviality, and composition laws the remainder of the monograph
interrogates: Chapter~\ref{ch05:main} proves that the legacy cover clauses
generate only trivial modalities, and the later constructive chapters build a
finite site on which two genuinely distinct nontrivial modalities exist and
commute.

Everything in this chapter is standard material. The definitions and theorems
are drawn from the textbook and reference literature
\cite{maclane,johnstone,stacks00W1,stacks00ZC}, and no novelty whatsoever is
claimed for them; they are cited as background throughout. What is specific to
this monograph is the selection and the emphasis: we dwell on exactly the
points---matching families on sieves containing identities, the basis
criterion, the behavior of sheafification on finite posets, and the criterion
for a modality to be trivial---that carry the weight of the collapse analysis
in Chapter~\ref{ch05:main} and of the finite constructions later. Examples are
chosen on small finite posets, deliberately miniature versions of the
sixteen-context lattice used by the successor construction, so that every
verification in this chapter can be redone by hand in a few minutes and every
later exhaustive machine check has a hand-checked prototype here.

\begin{notation}\label{ch04:not:conventions}
Throughout this chapter $\cC$ denotes a small category, $\PSh(\cC)$ the
category of presheaves $F:\cC^{\mathrm{op}}\to\Set$, and $yC=\Hom_{\cC}(-,C)$
the representable presheaf at $C\in\Ob(\cC)$. For $x\in F(C)$ and
$f:D\to C$ we write $x|_f$ for $F(f)(x)\in F(D)$, and when $\cC$ is a poset and
$U\le V$ we write $x|_U$ for the restriction of $x\in F(V)$ along the unique
arrow $U\to V$. Sieves, pullback sieves $g^{*}S$, the maximal sieve $m_C$, and
Grothendieck topologies $J$ are as in Chapter~3. A sieve on $C$ is identified
with a subfunctor of $yC$.
\end{notation}

\section{Sections, restriction, and the idea of local data}\label{ch04:sec:idea}

A presheaf is a system of ``data assignments'': to every object $C$ a set
$F(C)$ of \emph{sections over $C$}, and to every arrow $f:D\to C$ a
\emph{restriction} operation $F(C)\to F(D)$, contravariantly and
compatibly with composition. Nothing in the presheaf axioms relates sections
over an object to sections over the things that map into it, beyond the bare
existence of restriction. A presheaf may have many sections over $C$ that all
restrict to the same local data, or local data that assemble compatibly yet
come from no section over $C$ at all. Both phenomena will be exhibited on a
four-element poset in Section~\ref{ch04:sec:nonexamples}.

The sheaf condition is precisely the statement that neither pathology occurs
over covers. Informally: \emph{a section over $C$ is exactly a compatible
family of sections over any cover of $C$}. The slogan has two halves.
Uniqueness---two sections over $C$ agreeing on a cover are equal---is called
\emph{separation}. Existence---every compatible family of local sections glues
to a section over $C$---is the properly ``sheafy'' half. The definition below
makes ``compatible family'' precise for an arbitrary sieve; the version for a
generating family of arrows, which is what one actually checks in practice,
follows in Section~\ref{ch04:sec:basis}.

Why should an evidence-integration program care? The program's ambition,
recorded in the governing specification, is to treat locally gathered
observations as sections of presheaves over a site of contexts, and to let the
topology encode which families of contexts are jointly informative. Under that
reading, the sheaf condition is a \emph{composability discipline}: it says
which systems of local observations determine a global one and insists the
determination be unique. The mathematics of this chapter is entirely neutral
about any such interpretation---the theorems are about sets and arrows---but
the choice of what to prove is driven by it, and the claim boundaries of the
program (Ring~4, not licensed, no empirical interpretation) apply to every
interpretive gloss in this paragraph.

\section{Matching families and amalgamation}\label{ch04:sec:matching}

\begin{definition}[Matching family]\label{ch04:def:matching}
Let $F\in\PSh(\cC)$, let $C\in\Ob(\cC)$, and let $S$ be a sieve on $C$. A
\emph{matching family} for $S$ with values in $F$ is a function assigning to
every arrow $f:D\to C$ belonging to $S$ an element $x_f\in F(D)$, such that
for every arrow $h:E\to D$,
\[
x_{f\circ h} \;=\; x_f|_h \qquad\text{in } F(E).
\]
Note that $f\circ h\in S$ automatically, because $S$ is closed under
precomposition, so the condition is well posed. We write
$\mathrm{Match}(S,F)$ for the set of matching families.
\end{definition}

The compatibility equation quantifies over \emph{all} arrows $h$, not merely
over some chosen overlaps. This is the right level of generality for sieves;
when one works with a generating family and the category has pullbacks (or the
poset has meets), compatibility reduces to agreement on pairwise overlaps, as
we prove in Lemma~\ref{ch04:lem:meets} below.

\begin{remark}[Matching families are natural transformations]\label{ch04:rem:nat}
Regard the sieve $S$ as a subfunctor of $yC$: $S(D)\subseteq\Hom(D,C)$ is the
set of members of $S$ with domain $D$. Then a matching family for $S$ with
values in $F$ is exactly a natural transformation $S\Rightarrow F$: the
component at $D$ sends $f\in S(D)$ to $x_f$, and the compatibility equation is
exactly naturality. Hence
\[
\mathrm{Match}(S,F)\;=\;\Hom_{\PSh(\cC)}(S,F),
\]
and matching families inherit all the functoriality one expects: they can be
restricted along inclusions of sieves $S'\subseteq S$, pulled back along
arrows $g:D\to C$ (giving a matching family for $g^{*}S$), and pushed forward
along maps of presheaves $F\to G$. These three operations are used constantly
below and in the plus construction of Section~\ref{ch04:sec:plus}.
\end{remark}

\begin{definition}[Amalgamation; separated; sheaf]\label{ch04:def:sheaf}
Let $S$ be a sieve on $C$ and $(x_f)_{f\in S}$ a matching family for $S$ with
values in $F$. An \emph{amalgamation} of the family is an element $x\in F(C)$
with $x|_f=x_f$ for every $f\in S$. Every $x\in F(C)$ determines a canonical
matching family $(x|_f)_{f\in S}$; amalgamation asks for the converse. Define
the \emph{gluing map}
\[
\mathrm{glue}_{S,F}\;:\;F(C)\longrightarrow \mathrm{Match}(S,F),\qquad
x\longmapsto (x|_f)_{f\in S}.
\]
Given a Grothendieck topology $J$ on $\cC$, a presheaf $F$ is
\begin{enumerate}
\item \emph{separated} for $J$ if $\mathrm{glue}_{S,F}$ is injective for every
$C$ and every $S\in J(C)$;
\item a \emph{sheaf} for $J$ if $\mathrm{glue}_{S,F}$ is bijective for every
$C$ and every $S\in J(C)$.
\end{enumerate}
We write $\Sh(\cC,J)\subseteq\PSh(\cC)$ for the full subcategory of sheaves.
A topology is \emph{subcanonical} if every representable presheaf is a sheaf
\cite{maclane,stacks00WO}.
\end{definition}

Under the identification of Remark~\ref{ch04:rem:nat}, the gluing map is
composition with the inclusion $S\hookrightarrow yC$ followed by the Yoneda
identification $\Hom_{\PSh}(yC,F)\cong F(C)$; the sheaf condition says exactly
that the inclusion of every covering sieve into its representable becomes
invisible to $F$. This one-line reformulation is the
source of both the triviality criterion of
Proposition~\ref{ch04:prop:trivial} and the collapse diagnosis of
Chapter~\ref{ch05:main}.

\begin{example}[Maximal sieves are free]\label{ch04:ex:maximal}
For every presheaf $F$ and every object $C$, the gluing map for the maximal
sieve $m_C$ is a bijection. Indeed a matching family $(x_f)_{f\in m_C}$ is
determined by its value at the identity: $x_f=x_{\id_C\circ f}=x_{\id_C}|_f$,
so the assignment $(x_f)\mapsto x_{\id_C}$ is inverse to
$\mathrm{glue}_{m_C,F}$. Consequently axiom (J1)---every maximal sieve
covers---imposes no condition at all on presheaves. A topology whose covering
sieves are \emph{only} the maximal ones therefore has every presheaf as a
sheaf. This innocuous computation is the engine of the legacy collapse: a
cover clause that generates only maximal sieves is a cover clause that
constrains nothing (Proposition~\ref{ch04:prop:trivial} and
Chapter~\ref{ch05:main}).
\end{example}

\begin{lemma}[Compatibility via meets on a poset]\label{ch04:lem:meets}
Let $\cC$ be a poset with binary meets, $V\in\cC$, and let
$\{U_i\le V\}_{i\in I}$ be a family generating the sieve
$S=\{W \mid W\le U_i \text{ for some } i\}$. Then a matching family for $S$
with values in $F$ is the same data as a tuple $(x_i\in F(U_i))_{i\in I}$
satisfying
\[
x_i|_{U_i\wedge U_j} \;=\; x_j|_{U_i\wedge U_j}\qquad\text{for all } i,j\in I.
\]
(Standard; see \cite[Chap.~III]{maclane}.)
\end{lemma}

\begin{proof}
Given a matching family $(x_W)_{W\in S}$, set $x_i=x_{U_i}$; compatibility at
$U_i\wedge U_j\le U_i,U_j$ follows from the matching condition applied to the
two factorizations of $U_i\wedge U_j\to V$. Conversely, given a compatible
tuple, define $x_W=x_i|_W$ for any $i$ with $W\le U_i$; this is well defined
because if also $W\le U_j$ then $W\le U_i\wedge U_j$ and the two candidate
values agree by compatibility after restriction. The matching condition for
$(x_W)$ is then immediate, since restriction is functorial and every arrow of a
poset is a unique inequality. The two constructions are mutually inverse.
\end{proof}

Lemma~\ref{ch04:lem:meets} is why finite poset examples are so pleasant: a
matching family is just a compatible tuple, and both the matching-family set
and the section set are finite sets one can list. All the finite sites in this
monograph---the miniature examples below, and the sixteen-context lattice of
the successor---are posets with meets (intersections), so the lemma applies
verbatim throughout.

\section{First examples on a small poset}\label{ch04:sec:examples}

Fix the four-element poset
\[
L_2\;=\;\mathcal{P}(\{a,b\})\;=\;\{\varnothing,\ \{a\},\ \{b\},\ \{a,b\}\},
\]
ordered by inclusion. It has meets (intersection) and joins (union). Equip it
with the coverage whose generating data are: the two-element family
$\{\{a\}\to\{a,b\},\ \{b\}\to\{a,b\}\}$ covering the top; the identity family
at each object; and the \emph{null cover}---the empty family---covering
$\varnothing$. Write $J_2$ for the generated topology. (That these data
generate a topology is an instance of the basis machinery of
Section~\ref{ch04:sec:basis}; the reader may also verify (J1)--(J3) directly,
a worthwhile finger exercise.) The poset $L_2$ is the two-atom miniature of
the four-atom context lattice $L=\mathcal{P}(\{0,1\}\times\{0,1\})$ on which
the successor construction lives, and the covering pattern---partition the
top into its atoms, null-cover the bottom---is the miniature of the partition
coverages there.

Before any example we fix one point. The null cover at
$\varnothing$ is a genuine covering sieve (the empty sieve), and the sheaf
condition for it reads: $F(\varnothing)\to\mathrm{Match}(\varnothing\text{-sieve},F)$
is a bijection. A matching family for the empty sieve is a function defined on
no arrows; there is exactly one such function. Hence:

\begin{lemma}[The null cover forces a singleton]\label{ch04:lem:null}
If the empty sieve covers an object $Z$, then every sheaf $F$ has
$F(Z)$ a one-element set. (Standard; the two-line proof is self-contained.)
\end{lemma}

\begin{proof}
$\mathrm{Match}(\varnothing,F)$ is a singleton (the empty family, vacuously
matching), and the gluing map $F(Z)\to\mathrm{Match}(\varnothing,F)$ must be a
bijection.
\end{proof}

\begin{example}[The function presheaf is a sheaf]\label{ch04:ex:functions}
Fix a set $X$ and let $F(V)=\Set(V,X)$, the set of all functions from the
subset $V\subseteq\{a,b\}$ to $X$, with restriction of functions as
restriction. Then $F$ is a $J_2$-sheaf. At the null cover:
$F(\varnothing)=\Set(\varnothing,X)$ is a singleton (the empty function), as
Lemma~\ref{ch04:lem:null} requires. At the top cover: by
Lemma~\ref{ch04:lem:meets} a matching family is a pair
$(s\in\Set(\{a\},X),\ t\in\Set(\{b\},X))$ agreeing on
$\{a\}\cap\{b\}=\varnothing$---a vacuous condition since
$F(\varnothing)$ is a singleton---so
$\mathrm{Match}\cong X\times X$; and $F(\{a,b\})=\Set(\{a,b\},X)\cong
X\times X$ by evaluation at $a$ and $b$. The gluing map is exactly this
identification, hence bijective. Identity covers are free by
Example~\ref{ch04:ex:maximal}. This is the paradigm sheaf: sections are
literally functions, and functions are determined by, and glued from, their
values on the pieces of a partition.
\end{example}

\begin{example}[Representables; subcanonicity of $J_2$]\label{ch04:ex:representables}
For $W\in L_2$ the representable $yW$ has $yW(V)$ a singleton if $V\subseteq W$
and empty otherwise. Every representable is a $J_2$-sheaf, so $J_2$ is
subcanonical. The check at the top cover splits into cases. If
$W=\{a,b\}$: sections over $\{a\}$, $\{b\}$, and the top are all singletons,
matching families are singletons, gluing is a bijection of singletons. If
$W=\{a\}$: then $yW(\{b\})=\varnothing$, so there is \emph{no} matching family
at all for the top cover (a matching family must supply an element over every
member of the sieve, including $\{b\}\to\{a,b\}$), and also
$yW(\{a,b\})=\varnothing$; the gluing map is the unique map
$\varnothing\to\varnothing$, a bijection. The case $W=\{b\}$ is symmetric, and
$W=\varnothing$ likewise gives $\varnothing\to\varnothing$ at the top and
singletons at the null cover ($\varnothing\subseteq W$ always). The general
shape of this argument---if every piece of a partition of $V$ lies inside $W$
then $V$ itself lies inside $W$, because $V$ is the union of the
pieces---reappears as the subcanonicity theorem for the full context lattice
in the successor development, and stands in instructive contrast to
Chapter~\ref{ch06:main}, where a jointly-surjective coverage on a category
with \emph{too few} morphisms fails subcanonicity.
\end{example}

\begin{example}[A top-only side condition: separated but not a sheaf]\label{ch04:ex:bounded}
Let $X=\{0,1\}$ and let $G'(V)\subseteq\Set(V,X)$ consist of all functions
when $V\neq\{a,b\}$, and of those functions that are \emph{not} constantly
$1$ when $V=\{a,b\}$. Restriction preserves membership, since no condition is
imposed below the top, so $G'$ is a subpresheaf of the sheaf $\Set(-,X)$ of
Example~\ref{ch04:ex:functions} and is therefore separated (injectivity of
gluing is inherited by subpresheaves). It is not a sheaf: $s(a)=1$ and
$t(b)=1$ are both allowed locally, they match vacuously over
$\{a\}\cap\{b\}=\varnothing$, but any amalgamation would be the constant
function $1$ on $\{a,b\}$, which $G'$ excludes. This is the standard failure
mode of ``global side conditions'': a property imposed at the top but not
stable under gluing from pieces destroys sheafhood without harming
separation. Imposing the same condition at \emph{every} context would not
even define a presheaf---the restriction to $\{a\}$ of the function with
values $1,0$ on $\{a,b\}$ is constantly $1$---which is a reminder that
restriction-stability is the first thing to check.
\end{example}

\section{Non-examples, and what each failure teaches}\label{ch04:sec:nonexamples}

Non-examples are where the sheaf condition earns its keep. We record three on
$L_2$, each isolating one failure mode, and summarize the verdicts in
Table~\ref{ch04:tab:verdicts}.

\begin{example}[The constant presheaf fails the null cover]\label{ch04:ex:constant}
Fix a set $X$ with at least two elements and let $\Delta_X(V)=X$ for
\emph{every} $V$, all restrictions the identity. At the top cover,
matching families are pairs $(x,y)\in X\times X$ with $x|_\varnothing =
y|_\varnothing$, i.e.\ $x=y$ since restrictions are identities; so
$\mathrm{Match}\cong X$, and gluing $X\to X$ is the identity---fine. But at
the null cover, Lemma~\ref{ch04:lem:null} demands
$|\Delta_X(\varnothing)|=1$, whereas $|X|\ge 2$. So $\Delta_X$ is not a
$J_2$-sheaf, and the failure is entirely concentrated at the empty context.
The lesson: the empty cover is not a degenerate afterthought; it normalizes
the value at the initial context, and every finite construction later in the
monograph (the reflector formula in particular, whose empty product clause
sets the value at $\varnothing$ to a singleton) must and does handle it
explicitly.
\end{example}

\begin{example}[Separated but not a sheaf: missing gluings]\label{ch04:ex:separated}
Let $X$ have at least two elements. Define $H(\varnothing)=\{*\}$,
$H(\{a\})=H(\{b\})=H(\{a,b\})=X$, with the restrictions between nonempty
contexts the identity on $X$ and the restrictions to $\varnothing$ the unique
map to the singleton. At the top cover, a matching family is any pair
$(x,y)\in X\times X$: the compatibility condition lives in
$H(\varnothing)=\{*\}$ and is vacuous. An amalgamation of $(x,y)$ is an
element $z\in X$ with $z=x$ and $z=y$; it exists if and only if $x=y$ and is
then unique. So $\mathrm{glue}$ is injective but not surjective: $H$ is
separated, not a sheaf. The presheaf ``thinks'' a section over the top is a
single element of $X$, but the cover sees two independent local choices. Its
sheafification, computed in Example~\ref{ch04:ex:plusworked} below, repairs
this by \emph{creating} the missing sections: the associated sheaf has
$X\times X$ over the top. The witness presheaves that separate the two
modalities in the successor's finite theory are exactly presheaves of this
character: sheaf-like along one partition, gluing-deficient along the other.
\end{example}

\begin{example}[Not separated: redundant global sections]\label{ch04:ex:nonseparated}
Define $K(\{a,b\})=\{u,v\}$ with $u\ne v$, and $K(V)=\{*\}$ for all proper
$V\subseteq\{a,b\}$, with the only possible restrictions. At the top cover
there is exactly one matching family (all values are forced into singletons),
and \emph{both} $u$ and $v$ are amalgamations of it: gluing is surjective but
not injective. $K$ is not separated. The presheaf carries a distinction at the
top---$u$ versus $v$---that no cover can see; sheafification will identify
$u$ and $v$. In audit language: $K$ records a difference with no observational
support on the covering contexts, and passage to sheaves deletes exactly that
difference. Whether such deletion is desirable is a modeling decision, not a
mathematical one; the mathematics only says that \emph{if} one demands the
sheaf discipline, unsupported distinctions cannot survive.
\end{example}

\begin{table}[ht]
\centering
\begin{tabular}{lllll}
\toprule
Presheaf & Defined by & Null cover & Top cover & Verdict \\
\midrule
$F=\Set(-,X)$ & functions into $X$ & singleton & bijective & sheaf \\
$yW$ & representable & singleton & bijective & sheaf \\
$G'$ & top-only side condition & singleton & not surjective & separated only \\
$\Delta_X$, $|X|\ge2$ & constant & fails & bijective & not a sheaf \\
$H$, $|X|\ge2$ & identity restrictions & singleton & not surjective & separated only \\
$K$ & doubled top & singleton & not injective & not separated \\
\bottomrule
\end{tabular}
\caption{Verdicts for the presheaves of
Sections~\ref{ch04:sec:examples}--\ref{ch04:sec:nonexamples} on the poset
$L_2$ with topology $J_2$.}
\label{ch04:tab:verdicts}
\end{table}

\begin{exercise}\label{ch04:exr:census}
$L_2$ has four objects; list all sieves on each object (there are six on the
top, three on each singleton, two on $\varnothing$) and determine exactly
which collections of sieves form Grothendieck topologies on $L_2$. Verify
that $J_2$ is the smallest topology in which the atom family covers the top
and the empty sieve covers $\varnothing$. This is the two-atom warm-up for
the $298$-sieve census that the finite certificate performs on the
sixteen-context lattice.
\end{exercise}

\section{The basis criterion}\label{ch04:sec:basis}

Definition~\ref{ch04:def:sheaf} quantifies over all covering sieves of a
topology, which even on the sixteen-object successor lattice means every
covering sieve of the topology in question---$168$ of them for the joined
topology $J_A\vee J_R$ (Corollary~\ref{ch09:cor-census}), and a strictly
smaller collection for each of the coarser $J_A$ and $J_R$, not separately
enumerated in this monograph. In practice one never checks them all
directly; one checks a \emph{generating} family of covers and appeals to the
basis criterion. This
section states the criterion, proves it in the finite-poset generality used
throughout the monograph, and cites the general statement.

\begin{definition}[Coverage basis]\label{ch04:def:basis}
A \emph{coverage basis} (or pretopology-style basis) on $\cC$ assigns to each
object $C$ a set $K(C)$ of families of arrows with codomain $C$, such that:
\begin{enumerate}
\item (Identity) the singleton family $\{\id_C\}$ belongs to $K(C)$;
\item (Stability) for every $\{f_i:C_i\to C\}\in K(C)$ and every arrow
$g:D\to C$, there is a family $\{h_j:D_j\to D\}\in K(D)$ such that every
composite $g\circ h_j$ factors through some $f_i$;
\item (Composition) if $\{f_i:C_i\to C\}\in K(C)$ and for each $i$ a family
$\{g_{ik}:C_{ik}\to C_i\}\in K(C_i)$ is given, then the sieve generated by
the composite family $\{f_i\circ g_{ik}\}_{i,k}$ contains the sieve generated
by some family in $K(C)$.
\end{enumerate}
This is the coverage (refinement) form of the axioms;
Definition~\ref{ch03:def:pretopology} gave the stricter pretopology form,
which presupposes pullbacks and states stability as literal pullback of
families. The successor's partition families satisfy both forms; the $J_2$
data of Section~\ref{ch04:sec:examples} satisfy the coverage form only, since
the pullback $\{\{a\},\varnothing\}$ of the atom family is not itself a
declared generating family at $\{a\}$. The arguments below use only the
coverage form.

The \emph{generated topology} $J_K$ is the smallest Grothendieck topology in
which, for every $C$ and every family in $K(C)$, the sieve generated by that
family is covering. Smallest makes sense because the collection of topologies
containing a given set of sieves is closed under intersection; this closure
point is developed at length in Chapter~\ref{ch05:main}, where it becomes the
provenance-erasure argument. See \cite{stacks00ZC} and
\cite[Chap.~III]{maclane} for the standard treatment.
\end{definition}

\begin{theorem}[Basis criterion; general case cited, finite-poset case proved
here]\label{ch04:thm:basis}
Let $K$ be a coverage basis on $\cC$ and $J_K$ the generated topology. A
presheaf $F$ is a $J_K$-sheaf if and only if for every object $C$ and every
generating family $\{f_i:C_i\to C\}\in K(C)$, the gluing map from $F(C)$ to
matching families for the sieve generated by $\{f_i\}$ is a bijection. The
general statement, for an arbitrary small category with a coverage basis, is
standard background \cite{stacks00ZC,maclane} and is not reproved here. The
proof below establishes the case of finite posets with meets in which $K$ is
moreover \emph{strongly pullback-stable}: for every
$\{f_i:C_i\to C\}\in K(C)$ and every $U\le C$, the pulled-back family
generates a sieve on $U$ that is itself basis-generated or maximal. That is
the only case used in this monograph, and the hypothesis has to be verified
per basis rather than assumed.
\end{theorem}

The forward implication is trivial (generated sieves are covering). The
substantive converse is standard \cite{stacks00ZC,maclane}; we prove it in the
setting of finite posets with meets, which covers every site analyzed
computationally in this monograph and exposes the mechanism---saturation---in
its executable form.

\begin{algorithm}[Saturation of a basis on a finite poset]\label{ch04:alg:saturation}
Input: a finite poset $P$ with meets and a coverage basis $K$. Initialize
$\mathcal{S}$ as the set of all pairs $(V,S)$ where $S$ is either a maximal
sieve or the sieve generated by a member of $K(V)$. Repeat until no change:
\begin{enumerate}
\item (Maximality) for every $(V,S)\in\mathcal{S}$ with $S=m_V$, keep;
add $(V,m_V)$ for every $V$;
\item (Pullback) for every $(V,S)\in\mathcal{S}$ and every $U\le V$, add
$(U,\,U^{*}S)$ where $U^{*}S=\{W\le U \mid W\in S\}$;
\item (Local character) for every $V$ and sieve $T$ on $V$: if some
$(V,S)\in\mathcal{S}$ has, for every $W\in S$, the pair
$(W,\,W^{*}T)\in\mathcal{S}$, add $(V,T)$.
\end{enumerate}
Output the fixed point $\mathcal{S}^{\infty}$.
\end{algorithm}

\begin{lemma}\label{ch04:lem:saturation}
On a finite poset, Algorithm~\ref{ch04:alg:saturation} terminates and its
fixed point is exactly the generated topology $J_K$.
\end{lemma}

\begin{proof}
Termination: the set of (object, sieve) pairs is finite and $\mathcal{S}$
grows monotonically. The fixed point satisfies (J1) by step (1), (J2) by step
(2), and (J3) by step (3), so it is a topology containing the generated
sieves; conversely each step only adds sieves that every topology containing
the generators must contain---(J1), (J2), (J3) are exactly the three
closure rules---so by induction on the stage at which a pair enters
$\mathcal{S}$, the fixed point is contained in every such topology. Hence it
equals the smallest one.
\end{proof}

This is not merely an existence device: the successor's modality-composition
verifier literally runs this saturation on the sixteen-context lattice,
enumerating all $298$ sieves and reporting $168$ covering ones, and then
re-verifies (J1)--(J3) on the computed fixed point as an independent check.
Algorithm~\ref{ch04:alg:saturation} is the hand-sized description of that
computation.

\begin{proof}[Proof of Theorem~\ref{ch04:thm:basis} for finite posets with meets]
Suppose $F$ glues bijectively along every basis-generated sieve; we show $F$
glues bijectively along every $S\in J_K(V)$, by induction over the stages of
Algorithm~\ref{ch04:alg:saturation}. The induction hypothesis is carried in
pullback-closed form: what is proved for a pair $(V,S)$ is that $F$ glues
bijectively along $U^{*}S$ for \emph{every} $U\le V$, the case $U=V$
included. The strengthening is what discharges the pullback step, and it is
needed, because the strong stability hypothesis is a hypothesis about
generating families only: a sieve that enters at the local-character step
need not contain any basis-generated sieve, so that hypothesis says nothing
about its pullbacks. Stage $0$: basis-generated sieves hold by hypothesis;
maximal sieves hold by Example~\ref{ch04:ex:maximal}.

Stage $0$ in pullback-closed form: let $S$ be generated by a family
$\{V_{i}\to V\}\in K(V)$ and let $U\le V$. On a poset,
$U^{*}S=\{W\le U\mid W\in S\}$, and on a poset with meets this is precisely
the sieve generated by the pulled-back family $\{U\wedge V_{i}\to U\}$. By
the strong stability hypothesis of Theorem~\ref{ch04:thm:basis} that family
generates a basis-generated or maximal sieve at $U$, so $U^{*}S$ is again
basis-generated or maximal and the base case applies to it directly; and if
$S=m_{V}$ then $U^{*}S=m_{U}$.
That hypothesis must be checked per basis, not assumed of ``the poset
setting'': for the successor's partition bases it is
Lemma~\ref{ch09:lem-basis} (intersection with $U$ carries blocks to blocks);
for the $J_2$ basis of Section~\ref{ch04:sec:examples} the pullback of the
atom family along $\{a\}\subseteq\{a,b\}$ is $\{\{a\},\varnothing\}$, which
is \emph{not} a declared generating family at $\{a\}$, but it generates
$m_{\{a\}}$, so the hypothesis holds there as well. (This is also the point
where a general-category proof needs more care; the general statement is
\cite{stacks00ZC}.)

Pullback step: a pair $(U,U^{*}S)$ added at step (2) from $(V,S)$ needs no
argument of its own, since $U'^{*}(U^{*}S)=U'^{*}S$ for every $U'\le U$: the
pullback-closed statement for $(V,S)$ already contains the pullback-closed
statement for $(U,U^{*}S)$. This is the step at which the unstrengthened
induction would appeal to strong stability outside its scope.

Local-character step: suppose $T$ is a sieve on $V$, $S\in J_K(V)$ satisfies
the bijective-gluing conclusion (inductively, together with all its
pullbacks), and for every $W\in S$ the pullback $W^{*}T$ does too. Injectivity
of gluing along $T$: if $x,y\in F(V)$ agree on all of $T$, then for each
$W\in S$, the restrictions $x|_W,y|_W$ agree on $W^{*}T$, which glues
bijectively over $W$; hence $x|_W=y|_W$ for all $W\in S$, and gluing along
$S$ (injective) gives $x=y$. Surjectivity: let $(x_W)_{W\in T}$ be a matching
family for $T$. For each $W\in S$, the family restricts to a matching family
for $W^{*}T$, which glues to a unique $z_W\in F(W)$. The $z_W$ form a matching
family for $S$: for $W'\le W$ in $S$, both $z_W|_{W'}$ and $z_{W'}$ glue the
restricted family on $W'^{*}T$, and uniqueness of gluing along $W'^{*}T$
forces them equal. Gluing along $S$ yields $z\in F(V)$; one checks
$z|_{W}=x_{W}$ for $W\in T$ by comparing both sides on the covering sieve
$W^{*}S$ of $W$ and using injectivity there (each $W''\in W^{*}S$ lies in both
$S$-land and $T$-land, where the values agree by construction). Hence gluing
along $T$ is bijective. The pullback-closed form for $(V,T)$ follows by
rerunning the same argument at each $U\le V$ with $U^{*}S$ in place of $S$:
gluing along $U^{*}S$ and along all of its pullbacks is available from the
pullback-closed hypothesis for $(V,S)$, and each $W\in U^{*}S$ lies in $S$
with $W^{*}(U^{*}T)=W^{*}T$, so the pullback-closed hypothesis for
$(W,W^{*}T)$ supplies the rest. This completes the induction.
\end{proof}

\begin{remark}\label{ch04:rem:basispractice}
The practical import: to certify sheafhood on the successor lattice one checks
each presheaf against the generating partition families only---two families
per nonempty context plus the null cover---rather than against every covering
sieve of the topology concerned ($168$ for the join $J_A\vee J_R$, fewer for
each of $J_A$ and $J_R$). The finite certificate does exactly
this for the sixteen representables and the two witness presheaves, and
Theorem~\ref{ch04:thm:basis} is the license. Conversely, when
\emph{refuting} sheafhood (as for the witness presheaves against the
transverse topology, or for the representable in the countermodel of
Chapter~\ref{ch06:main}), a single failed generating family suffices, since
generated sieves are covering.
\end{remark}

\section{The equalizer formulation, and life beyond posets}\label{ch04:sec:equalizer}

The matching-family formulation of the sheaf condition is the right one for
sieves; for a generating family of arrows there is an older and equally
useful formulation as an equalizer, which we record both for completeness and
because it is the form in which gluing computations are usually organized in
the literature \cite{maclane,stacks00ZC}.

\begin{proposition}[Equalizer form]\label{ch04:prop:equalizer}
Let $\cC$ have pullbacks, let $\{f_i:C_i\to C\}_{i\in I}$ be a family of
arrows generating the sieve $S$, and let $F\in\PSh(\cC)$. Then the gluing map
for $S$ is a bijection if and only if the diagram
\[
F(C)\;\longrightarrow\;\prod_{i\in I}F(C_i)
\;\rightrightarrows\;\prod_{i,j\in I}F(C_i\times_C C_j)
\]
is an equalizer of sets, where the first map is restriction along the $f_i$
and the parallel pair restricts along the two pullback projections
$C_i\times_C C_j\to C_i$ and $C_i\times_C C_j\to C_j$.
\end{proposition}

\begin{proof}[Proof sketch]
An element of the equalizer is a tuple $(x_i)$ with
$x_i|_{\pi_1}=x_j|_{\pi_2}$ on every pullback. Given such a tuple, define
$x_f=x_i|_g$ for any factorization $f=f_i\circ g$ of a member $f\in S$; the
pullback compatibility makes this independent of the chosen factorization
(two factorizations through $f_i$ and $f_j$ induce a map to
$C_i\times_C C_j$ through which both restrictions agree), and the resulting
assignment is a matching family. Conversely a matching family restricts to an
equalizer element on the generators. The two constructions are mutually
inverse; details are in \cite[Chap.~III]{maclane}.
\end{proof}

On a poset with meets, $C_i\times_C C_j$ is the meet $U_i\wedge U_j$ and the
proposition specializes to Lemma~\ref{ch04:lem:meets}. Two warnings attach to
the general case, and both matter later. First, if $\cC$ lacks the relevant
pullbacks, the equalizer form is not available and one must either use the
sieve form directly or check compatibility over \emph{all} pairs of
factorizations, as in Definition~\ref{ch04:def:matching}; the sieve form is
the primitive notion, the equalizer form a convenience. Second, even when
pullbacks exist, a family can have empty or degenerate pullbacks that make
the compatibility condition vacuous while gluing still fails---this is
exactly the geometry of the countermodel in Chapter~\ref{ch06:main}, where
two point inclusions into a two-element set have empty overlap, so every pair
of local sections is ``compatible,'' yet no global section exists. Vacuous
compatibility plus failed gluing is the signature of a cover that is
surjective on points but not effective on morphisms, and the reader should
hold the $L_2$ intuition (Example~\ref{ch04:ex:functions}) and that failure
side by side: the function presheaf glues because functions out of a union
are determined by functions on the pieces; a representable in an
injection-only category may not glue because the would-be glued map is not an
arrow of the category at all.

\begin{example}[A non-poset site in miniature]\label{ch04:ex:nonposet}
Let $\cC$ be the category with two objects $D$, $C$, two parallel arrows
$u,v:D\to C$, and identities. Let $S$ be the sieve on $C$ generated by
$\{u,v\}$; its members are $u$, $v$ (and nothing else, as there is nothing to
precompose with except identities). A matching family for $S$ valued in $F$
is a pair $(x_u,x_v)$ with $x_u,x_v\in F(D)$ and \emph{no} compatibility
condition, since there are no nonidentity arrows into $D$. The gluing map
sends $x\in F(C)$ to $(x|_u,x|_v)$. For $F=yD$: $yD(C)=\varnothing$ (no
arrows $C\to D$) while $yD(D)=\{\id_D\}$, so the matching-family set is a
singleton and the gluing map $\varnothing\to\{*\}$ is not surjective---the
representable $yD$ fails gluing for any topology declaring $S$ covering. This
tiny computation is the structural germ of the subcanonicity countermodel of
Chapter~\ref{ch06:main}: nothing about it is exotic, and its entire content
is that the category supplies no arrow to serve as the glued section. It also
illustrates why Definition~\ref{ch04:def:matching} quantifies over all
precompositions: on posets that quantification collapses to meets, but in
general the sieve's own arrow structure is the compatibility data.
\end{example}

\begin{remark}[Separated presheaves as an intermediate reflection]
\label{ch04:rem:separated}
Between $\PSh(\cC)$ and $\Sh(\cC,J)$ sits the full subcategory of separated
presheaves. It too is reflective: the first plus construction (Section
\ref{ch04:sec:plus}) composed with an image factorization reflects a presheaf
onto its separated quotient, identifying sections that agree on some cover.
The two-step structure of sheafification---first kill unsupported
distinctions, then adjoin missing gluings---is thus a composite of two
reflections, and the two steps are conceptually independent: Example
\ref{ch04:ex:nonseparated} needs only the first, Example
\ref{ch04:ex:separated} only the second. In audit terms the two steps have
different characters: separation \emph{deletes} data (a quotient), gluing
\emph{creates} data (new sections from compatible families), and a
provenance-conscious pipeline must track both, since neither is injective on
information. This is one of the structural reasons the successor stores
provenance outside the topology; the inside-the-topos operations are
irreversibly lossy by design.
\end{remark}

\begin{remark}[The canonical topology]\label{ch04:rem:canonical}
For every small category $\cC$ there is a largest subcanonical topology, the
\emph{canonical topology}, whose covering sieves are the \emph{universally
effective-epimorphic} sieves: those $S$ on $C$ such that every representable
glues uniquely along $S$ and along every pullback $g^{*}S$, $g:D\to C$
\cite{maclane,johnstone}. Subcanonicity of a given topology $J$ is exactly
the containment of $J$ in the canonical topology. This gives the right frame
for two later results: the successor's partition topologies on the context
lattice are subcanonical (they sit inside the canonical topology, as the
representable checks show), while the legacy jointly-surjective coverage on
the countermodel category of Chapter~\ref{ch06:main} is \emph{not} contained
in the canonical topology of that category---joint surjectivity on points
simply does not imply effective epimorphy of the family. ``Subcanonical'' is
thus not a decoration but a containment statement with content, and it can
fail.
\end{remark}

\section{Sheafification: the universal property}\label{ch04:sec:sheafification}

Not every presheaf is a sheaf, but every presheaf has a best sheaf
approximation, characterized by a universal property, constructed by the plus
construction, and enjoying left exactness. We take these in order.

\begin{definition}[Associated sheaf]\label{ch04:def:assoc}
Let $(\cC,J)$ be a site and $F\in\PSh(\cC)$. An \emph{associated sheaf} (or
\emph{sheafification}) of $F$ is a sheaf $aF$ together with a map
$\eta_F:F\to aF$ such that for every sheaf $G$, composition with $\eta_F$ is a
bijection
\[
\Hom_{\PSh}(aF,\,G)\;\xrightarrow{\ \cong\ }\;\Hom_{\PSh}(F,\,G).
\]
\end{definition}

\begin{proposition}\label{ch04:prop:uniqueness}
An associated sheaf, if it exists, is unique up to unique isomorphism
compatible with the units; and the assignment $F\mapsto aF$ extends uniquely
to a functor $a:\PSh(\cC)\to\Sh(\cC,J)$ left adjoint to the inclusion
$i:\Sh(\cC,J)\hookrightarrow\PSh(\cC)$.
\end{proposition}

\begin{proof}
Both statements are formal consequences of the universal property, by the
standard argument for reflections \cite{maclanecwm,riehl}: uniqueness because
two reflections of the same object corepresent the same functor on the
subcategory; functoriality because $\eta_{F'}\circ\alpha:F\to aF'$ extends
uniquely across $\eta_F$ for any $\alpha:F\to F'$; adjointness because the
defining bijection is exactly the adjunction isomorphism, and its naturality
follows from uniqueness of extensions.
\end{proof}

\begin{theorem}[Existence and exactness; background]\label{ch04:thm:sheafification}
For every small site $(\cC,J)$ the associated-sheaf functor
$a:\PSh(\cC)\to\Sh(\cC,J)$ exists, and it is \emph{left exact}: it preserves
all finite limits. Consequently $\Sh(\cC,J)$ is a reflective subcategory of
$\PSh(\cC)$ whose reflector preserves finite limits, and $\Sh(\cC,J)$ is a
Grothendieck topos.
\end{theorem}

This is textbook material \cite[Chap.~III, \S\S5--6 and Chap.~V, \S3]{maclane},
\cite{stacks00W1}, \cite{johnstone}; we present the construction in the next
section both because the monograph's finite computations imitate it and
because seeing it makes the left-exactness statement believable rather than
oracular. Left exactness matters for the program for a concrete reason: a
left-exact reflector preserves monomorphisms, finite products, and equalizers,
hence preserves the shape of relational and subobject structure when evidence
objects are pushed from the presheaf level into sheaves. A right adjoint that
merely existed, without exactness, would permit the reflection to distort
exactly the finite-limit structure that admissibility relations and typed
comparisons live in.

\section{The plus construction, and computing on finite sites}\label{ch04:sec:plus}

\begin{construction}[The plus construction]\label{ch04:con:plus}
For $F\in\PSh(\cC)$ and $C\in\Ob(\cC)$, the covering sieves of $C$ form a
directed poset under \emph{refinement}: $S'$ refines $S$ when $S'\subseteq S$,
and any two covering sieves have a common covering refinement
$S\cap S'$ (an instance of (J2) plus (J3)). Restriction of matching families
along inclusions gives a directed system, and one defines
\[
F^{+}(C)\;=\;\operatorname*{colim}_{S\in J(C)}\ \mathrm{Match}(S,F),
\]
the colimit over ever-finer covering sieves. An element of $F^{+}(C)$ is an
equivalence class of matching families, two families being identified when
they agree on a common covering refinement. Restriction along $g:D\to C$ sends
the class of $(S,(x_f))$ to the class of $(g^{*}S,(x_{g\circ h})_h)$; the unit
$F\to F^{+}$ sends $x\in F(C)$ to the class of the canonical family
$(x|_f)_{f\in m_C}$.
\end{construction}

\begin{theorem}[Background; \cite{stacks00W1,maclane}]\label{ch04:thm:plus}
$F^{+}$ is a separated presheaf for every $F$; if $F$ is separated, $F^{+}$ is
a sheaf. Hence $aF=F^{++}$, with unit the composite $F\to F^{+}\to F^{++}$,
satisfies the universal property of Definition~\ref{ch04:def:assoc}. Moreover
$a$ is left exact: each application of $(-)^{+}$ is a filtered colimit of
finite-limit-preserving operations $\mathrm{Match}(S,-)$, and filtered
colimits commute with finite limits in $\Set$.
\end{theorem}

Two applications of $(-)^{+}$ are needed in general: the first pass merges
sections that a cover cannot distinguish (repairing separation), the second
pass adjoins gluings for compatible families (repairing existence). On
particular presheaves one pass may happen to suffice, but no general theorem
says so, and the successor's proof artifacts are deliberately structured so as
never to assume it.

\begin{example}[A worked sheafification on $L_2$]\label{ch04:ex:plusworked}
Take $H$ from Example~\ref{ch04:ex:separated} ($H$ separated, identity
restrictions between nonempty contexts, $|X|\ge 2$). We compute $H^{+}$ for
$J_2$. On $L_2$ every object has a \emph{minimum} covering sieve: the
intersection of all covering sieves is covering (finitely many sieves, and
covering sieves are closed under pairwise intersection), so the directed
colimit collapses to its terminal stage:
\[
F^{+}(V)\;\cong\;\mathrm{Match}(S_{\min}(V),\,F)
\qquad\text{on any finite site.}
\]
Here $S_{\min}(\{a,b\})$ is the sieve generated by the two atoms,
$S_{\min}$ at each singleton is the maximal sieve, and
$S_{\min}(\varnothing)$ is the empty sieve. Therefore
$H^{+}(\{a,b\})=\mathrm{Match}(S_{\min},H)\cong X\times X$ (Example
\ref{ch04:ex:separated}), $H^{+}$ agrees with $H$ on singletons, and
$H^{+}(\varnothing)$ is a singleton. One checks directly that $H^{+}$ is
already a $J_2$-sheaf (it is isomorphic to the function presheaf
$\Set(-,X)$ of Example~\ref{ch04:ex:functions}), so here a single pass
sufficed---because $H$ was separated to begin with, exactly as
Theorem~\ref{ch04:thm:plus} predicts. Redoing the computation for the
non-separated $K$ of Example~\ref{ch04:ex:nonseparated} shows the first pass
collapsing $u$ and $v$ to a single class, after which the second pass has
nothing left to add.
\end{example}

\begin{remark}[Why direct constructions can beat the plus construction in
finite settings]\label{ch04:rem:direct}
On a finite site the plus construction is entirely computable, as
Example~\ref{ch04:ex:plusworked} shows, but it is rarely the best way to
\emph{certify} a sheafification, for three reasons. First, bookkeeping: the
colimit ranges over all covering sieves of the topology in question ($168$ of
them for the joined topology on the successor
lattice), and although the minimum-sieve collapse tames this, an implementation
still owes a proof that the collapse is legitimate on the exact site at hand.
Second, the one-pass question: whether $F^{+}$ is already a sheaf depends on
separation hypotheses that must be re-established for every input, and an
implementation that silently assumes one pass is wrong on non-separated
inputs. Third, and decisively for audit: the universal property, not the
construction, is the specification. A closed-form candidate reflector---such
as the product-over-blocks formula the successor uses on its partition
coverages, which assigns to each context the product of the presheaf's values
on the pieces---can be verified \emph{directly} against
Definition~\ref{ch04:def:assoc}: prove the candidate lands in sheaves, prove
composition with the unit is a bijection onto $\Hom_{\PSh}(F,G)$ for every
sheaf $G$, and conclude by uniqueness (Proposition~\ref{ch04:prop:uniqueness})
that the candidate \emph{is} sheafification, with no plus-construction
bookkeeping anywhere in the trusted path. The successor's
modality-composition proof takes exactly this route, and its stated reason is
the one just given: the direct proof avoids assuming that a single generic
plus construction sheafifies an arbitrary presheaf. In finite settings the
direct route is simultaneously more honest and easier to check exhaustively:
a formula evaluated at sixteen contexts, against a universal property checked
on finitely many fixtures, is a better certificate substrate than a two-pass
colimit whose intermediate objects are themselves quotients of function
spaces.
\end{remark}

\begin{exercise}\label{ch04:exr:plus}
Compute $\Delta_X^{+}$ and $\Delta_X^{++}$ on $(L_2,J_2)$ for $|X|=2$. Check
that $\Delta_X^{+}$ has one section over $\varnothing$, $X$ over each atom,
and only $X$ over the top---the matching condition through
$\Delta_X(\varnothing)$ forces the two atom values to agree---so that
$\Delta_X^{+}$ is separated but not yet a sheaf; then check that the second
pass enlarges the top to $X\times X$, so that
$a\Delta_X=\Delta_X^{++}\cong\Set(-,X)$, the function presheaf of
Example~\ref{ch04:ex:functions}. Say which pass deletes sections and which
creates them, and where each happens.
\end{exercise}

\section{Reflective subcategories and idempotent monads}\label{ch04:sec:reflective}

Sheafification exhibits $\Sh(\cC,J)$ as a reflective subcategory of
$\PSh(\cC)$. This section isolates the abstract shape, because the monograph's
later chapters manipulate several reflections at once and their composites,
and the abstract calculus is what makes those manipulations safe.

\begin{definition}\label{ch04:def:reflective}
A full subcategory $\mathcal{B}\subseteq\mathcal{A}$ is \emph{reflective} if
the inclusion $i:\mathcal{B}\to\mathcal{A}$ has a left adjoint
$a:\mathcal{A}\to\mathcal{B}$, the \emph{reflector}. The composite
$M=ia:\mathcal{A}\to\mathcal{A}$, together with the unit
$\eta:\mathrm{Id}\Rightarrow M$ and multiplication induced by the adjunction
counit, is a monad on $\mathcal{A}$ \cite{maclanecwm,riehl,awodey}.
\end{definition}

\begin{proposition}[Idempotency]\label{ch04:prop:idempotent}
For a reflective subcategory as above, the monad $M=ia$ is \emph{idempotent}:
the multiplication $\mu:MM\Rightarrow M$ is a natural isomorphism;
equivalently, $M\eta_F$ and $\eta_{MF}$ are equal isomorphisms
$MF\to MMF$ for every $F$.
\end{proposition}

\begin{proof}
Because $\mathcal{B}$ is full, the counit of the adjunction is an isomorphism
$a i\cong\mathrm{Id}_{\mathcal{B}}$ (the reflection of an object already in
$\mathcal{B}$ is itself, by the universal property applied to the identity).
Whiskering, $MM=iaia\cong ia=M$, and the triangle identities identify this
isomorphism with $\mu$; the two unit maps into $MMF$ agree because both are
the unique extension of $\eta_F$ across itself.
\end{proof}

Idempotency is what lets one speak of the subcategory and the monad
interchangeably: $\mathcal{B}$ is recovered inside $\mathcal{A}$ as the
objects for which the unit is an isomorphism, and $M$ is determined up to
canonical isomorphism by its image. The next definition packages the objects
side.

\begin{definition}[Local objects]\label{ch04:def:local}
Let $M=ia$ arise from a reflective subcategory. An object $F\in\mathcal{A}$ is
\emph{$M$-local} (or just \emph{local} when $M$ is understood) if
$\eta_F:F\to MF$ is an isomorphism.
\end{definition}

\begin{proposition}\label{ch04:prop:localssheaves}
For $M=ia$ with $a$ sheafification for a topology $J$ on $\cC$, an object of
$\PSh(\cC)$ is $M$-local if and only if it is a $J$-sheaf. In general, the
$M$-local objects are exactly the objects isomorphic to objects of the
reflective subcategory, and $M$ restricted to local objects is naturally
isomorphic to the identity.
\end{proposition}

\begin{proof}
If $F$ is a sheaf then $\eta_F$ is an isomorphism by the universal property
(extend $\id_F$ across $\eta_F$; the two composites are identities by
uniqueness of extensions). Conversely if $\eta_F$ is an isomorphism then $F$
is isomorphic to the sheaf $iaF$, and the sheaf condition---bijectivity of a
class of gluing maps---transfers across isomorphism. The general statements
are proved the same way, using only fullness and the universal property.
\end{proof}

\section{Modalities}\label{ch04:sec:modalities}

\begin{definition}[Modality]\label{ch04:def:modality}
A \emph{modality} on a presheaf category $\PSh(\cC)$ is a left-exact
idempotent monad $M$ on $\PSh(\cC)$ arising as $M=ia$ for a reflective
subcategory whose reflector $a$ preserves finite limits. A modality is
\emph{trivial} if it is naturally isomorphic to the identity functor. Two
modalities are \emph{distinct} if they are not naturally isomorphic;
equivalently (by Proposition~\ref{ch04:prop:localssheaves}) if their
subcategories of local objects differ.
\end{definition}

The word ``modality'' is borrowed deliberately. In topos theory, left-exact
reflective subcategories of a presheaf topos correspond to Lawvere--Tierney
topologies on it, which for $\PSh(\cC)$ correspond in turn to Grothendieck
topologies on $\cC$; the closure operators $j$ on the subobject classifier
that encode them behave like modal operators on the internal logic (``it is
locally the case that''), and the sheaf condition says exactly that a
predicate or object is insensitive to the modal blur. See
\cite[Chap.~V]{maclane} and \cite{johnstone} for this correspondence; in the
formal-logic and type-theoretic literature the same abstract shape---a
left-exact idempotent monad and its local objects---also travels under the
name modality. This monograph uses the term exclusively in the sense of
Definition~\ref{ch04:def:modality}, and always for modalities of
sheafification type: $M_J=i_Ja_J$ for a Grothendieck topology $J$ on the site
in question. No claim is made here about modal logics, and none is needed:
every later theorem about ``modalities'' is a theorem about sheafification
monads.

The program-level motivation, stated once and bounded: the governing
specification intends distinct modalities to serve as distinct
\emph{resolution disciplines} on evidence presheaves---one reflector
enforcing compatibility along one family of covers, another along a
transverse family---so that partially resolved objects (local for one
modality but not the other) are representable states of an audit pipeline.
For that intent to be realizable at all, the modalities in play must be
provably distinct and provably nontrivial, and their composites must be
understood. Those are precisely the questions the finite theory answers, and
the questions the legacy site fails, which is why the criterion below is the
hinge of the whole monograph.

\begin{proposition}[Triviality criterion]\label{ch04:prop:trivial}
Let $J$ be a Grothendieck topology on a small category $\cC$ and
$M_J=i_Ja_J$ the associated modality on $\PSh(\cC)$. The following are
equivalent:
\begin{enumerate}
\item $M_J$ is trivial (naturally isomorphic to the identity);
\item every presheaf on $\cC$ is a $J$-sheaf;
\item every covering sieve of $J$ is maximal: $J(C)=\{m_C\}$ for all $C$.
\end{enumerate}
This equivalence is standard; see \cite[Chap.~V]{maclane} and
\cite{johnstone}. The proof is given in full because
Chapter~\ref{ch05:main} applies the criterion directly to the legacy cover
clauses.
\end{proposition}

\begin{proof}
(3)$\Rightarrow$(2): Example~\ref{ch04:ex:maximal}---gluing along a maximal
sieve is bijective for every presheaf. (2)$\Rightarrow$(1): if every presheaf
is a sheaf then every unit $\eta_F:F\to M_JF$ is an isomorphism
(Proposition~\ref{ch04:prop:localssheaves}), and these assemble to a natural
isomorphism $\mathrm{Id}\cong M_J$. (1)$\Rightarrow$(3): we prove the
contrapositive with an explicit witness. Suppose some $S_0\in J(C)$ is not
maximal. Consider the presheaf of sieves
\[
\Omega(D)\;=\;\{\,\text{sieves on } D\,\},\qquad
\Omega(g:E\to D)\;=\;g^{*}(-),
\]
the subobject classifier of $\PSh(\cC)$; functoriality is the identity
$(g\circ h)^{*}=h^{*}g^{*}$ on sieves. Define a matching family for $S_0$
with values in $\Omega$ by assigning to every $f:D\to C$ in $S_0$ the maximal
sieve $m_D\in\Omega(D)$. This matches: $h^{*}m_D=m_E$ for every $h:E\to D$.
An amalgamation is a sieve $T$ on $C$ with $f^{*}T=m_D$ for all $f\in S_0$;
since $f^{*}T=m_D$ if and only if $\id_D\in f^{*}T$ if and only if $f\in T$,
the amalgamations are exactly the sieves $T\supseteq S_0$. There are at least
two: $S_0$ itself and $m_C$, distinct because $S_0$ is not maximal. So
$\Omega$ is not even separated for $J$, hence not a sheaf; if $M_J$ were
naturally isomorphic to the identity, $\Omega\cong M_J\Omega$ would be
isomorphic to a sheaf and hence a sheaf, a contradiction.
\end{proof}

\begin{remark}\label{ch04:rem:trivialreading}
Return to the proof of (1)$\Rightarrow$(3), because it is the
cleanest possible statement of what a nonmaximal cover \emph{does}: it forces
the presheaf of sieves to carry unsupported global distinctions (many sieves
over $C$ restrict to the maximal sieve everywhere on the cover), and
sheafification must therefore act nontrivially on $\Omega$. Conversely, a
topology all of whose covers are maximal never forces anything. The relevance
to the program is immediate and is spelled out in
Chapter~\ref{ch05:main}: the legacy ACBE and RPM clauses each generate only
maximal sieves, so by Proposition~\ref{ch04:prop:trivial} the modalities
built from either clause alone are trivial---naturally isomorphic to the
identity, local on every presheaf, incapable of supporting any distinction
between resolution states. The collapse theorem is, in this exact sense, an
instance of the triviality criterion, and the successor's entire design
problem is to manufacture covers that are honestly nonmaximal.
\end{remark}

\subsection{Closed sieves and the classifier viewpoint}\label{ch04:sub:closed}

The subobject classifier $\Omega$ of $\PSh(\cC)$, which appeared as the
witness in Proposition~\ref{ch04:prop:trivial}, supports a second description
of modalities, recorded here even though the monograph's proofs never
depend on it. Given a topology $J$, call a sieve $T$ on $C$
\emph{$J$-closed} if every arrow that lies in $T$ locally already lies in
$T$: whenever $f:D\to C$ has $f^{*}T\in J(D)$, then $f\in T$. Closing a sieve
(adding all such arrows) defines an operator $j:\Omega\to\Omega$, natural in
$C$, which is inflationary, idempotent, and meet-preserving; such an operator
on the classifier is a \emph{Lawvere--Tierney topology}, and on a presheaf
topos these correspond exactly to Grothendieck topologies on the base
category \cite[Chap.~V]{maclane}, \cite{johnstone}. The $J$-sheaves are then
the objects orthogonal to the class of \emph{dense} monomorphisms (those
whose closure is everything), and sheafification is the universal way of
inverting dense monos. Three consequences of this viewpoint are used later,
and we state them as background \cite{maclane,johnstone}:
\begin{enumerate}
\item The covering sieves of $J$ are exactly the sieves whose closure is
maximal; in particular a topology is determined by its closed sieves, and the
presheaf $\Omega_j$ of closed sieves is the subobject classifier of
$\Sh(\cC,J)$.
\item A modality of sheafification type inverts the inclusions
$S\hookrightarrow yC$ for covering $S$; this is the precise sense in which
``the topology cannot see the difference between a covering sieve and the
whole representable,'' and it is the abstract root of the provenance-erasure
discussion in Chapter~\ref{ch05:main}.
\item Meets of topologies correspond to pointwise meets of the associated
closure operators, and the join of two topologies is the topology generated
by their union---the operation the successor computes explicitly when it
identifies the join of its two partition topologies with the
singleton-partition topology.
\end{enumerate}
For the finite sites of this monograph, every one of these statements is also
directly checkable by enumeration, and the certificate chapters do so; the
classifier viewpoint is offered as orientation, not as a load-bearing
dependency.

\begin{exercise}\label{ch04:exr:closed}
On $(L_2,J_2)$, compute the closure of each of the six sieves on the top
object and verify: (i) the closed sieves on the top are exactly the maximal
sieve and the sieves generated by $\{a\}$, by $\{b\}$, and by the bottom
arrow $\varnothing\to\{a,b\}$ (note the \emph{empty} sieve is not closed,
since the null cover makes $\varnothing$ covered and so forces $\varnothing$
into every closure);
(ii) the covering sieves are exactly those whose closure is maximal;
(iii) $\Omega_j$ at the top has exactly four elements, matching the sheaf
condition's verdict on Exercise~\ref{ch04:exr:omega}.
\end{exercise}

\begin{remark}[Distinctness, composition, and what is \emph{not} claimed]
\label{ch04:rem:distinctness}
Two further points of the abstract calculus recur later. First, distinctness
of modalities is witnessed by objects: $M_J\ne M_{J'}$ (not naturally
isomorphic) as soon as some presheaf is local for one and not the other, since
reflectors fix exactly their local objects. This is how the successor's two
finite modalities are separated: an explicit presheaf that satisfies one
partition's gluing and fails the other's. Second, composites of modalities
are subtle in general---the composite of two left-exact idempotent monads
need not be idempotent, and when it is, identifying it requires a theorem,
not a reflex. On the successor's finite lattice the composite of the two
partition modalities is, by a proved and exhaustively corroborated canonical
natural isomorphism, the modality of the joined topology, and in particular
the two modalities \emph{commute}; distinctness does not imply
noncommutation, and no artifact of the program licenses a noncommutation
claim. Both facts are proved in the finite-theory chapters; they are
previewed here only so that the reader does not import an expectation of
order-sensitivity that the mathematics does not support.
\end{remark}

\begin{boundary}\label{ch04:bnd:status}
Status discipline for this chapter: every numbered statement above is either
(i) standard background, cited to \cite{maclane}, \cite{maclanecwm},
\cite{johnstone}, \cite{stacks00W1}, \cite{stacks00ZC}, \cite{riehl}, or
\cite{awodey}, with no novelty claimed; or (ii) an elementary finite
verification (the $L_2$ examples, Algorithm~\ref{ch04:alg:saturation} and its
lemma, and the finite-poset case of Theorem~\ref{ch04:thm:basis}) offered at
tutorial level. Forward references to properties of the successor
construction (distinct nontrivial commuting modalities, the $298$-sieve and
$168$-cover census, the direct reflector) describe results whose registered
status is \textsc{informal proof present} with exhaustive executable
corroboration; they are not peer-reviewed theorems, and nothing in this
chapter promotes them. No statement in this chapter carries any empirical,
causal, safety, or product interpretation.
\end{boundary}

\section{Chapter summary and the road ahead}\label{ch04:sec:summary}

The chapter assembled the four-piece toolkit the rest of the monograph runs
on. (1) \emph{The sheaf condition}: sections over an object biject with
matching families over each covering sieve; separation is the injective half,
gluing the surjective half; on posets with meets, matching families are
compatible tuples (Lemma~\ref{ch04:lem:meets}); the null cover normalizes the
initial context (Lemma~\ref{ch04:lem:null}). The examples and non-examples of
Sections~\ref{ch04:sec:examples}--\ref{ch04:sec:nonexamples} exhibit each
failure mode in the smallest possible habitat. (2) \emph{The basis
criterion}: sheafhood needs checking only on generating families
(Theorem~\ref{ch04:thm:basis}), and on finite posets the generated topology
is the fixed point of an explicit saturation
(Algorithm~\ref{ch04:alg:saturation})---the exact computation the finite
certificate later performs. (3) \emph{Sheafification}: a left-exact
reflection characterized by its universal property
(Definition~\ref{ch04:def:assoc}, Theorem~\ref{ch04:thm:sheafification}),
constructible by the two-pass plus construction
(Construction~\ref{ch04:con:plus}) but, on finite sites, better certified by
closed-form candidates verified directly against the universal property
(Remark~\ref{ch04:rem:direct}). (4) \emph{Modalities}: left-exact idempotent
monads of sheafification type, with local objects the sheaves
(Proposition~\ref{ch04:prop:localssheaves}), and with a sharp triviality
criterion---trivial exactly when every cover is maximal
(Proposition~\ref{ch04:prop:trivial}).

The criterion is the pivot on which the monograph now turns. Chapter
\ref{ch05:main} takes the registered legacy site, whose ACBE and RPM cover
clauses were intended to seed two distinct modalities, and shows by two
one-line sieve computations that every sieve either clause generates is
maximal. By Proposition~\ref{ch04:prop:trivial}, the intended modalities are
trivial; by the closure-operator analysis there, no relabeling, sub-sieve
surgery, or object-adjunction can help; and by
Chapter~\ref{ch06:main}, the same axioms also fail to deliver subcanonicity,
so even the representables must be handled through their associated sheaves.
Only after both negative results are fully absorbed does the successor
construction begin.

\begin{exercise}\label{ch04:exr:omega}
Compute $\Omega$ explicitly on $L_2$ (six sieves on the top, three on each
singleton, two on the bottom; compare Exercise~\ref{ch04:exr:census}),
identify the matching family used in the proof of
Proposition~\ref{ch04:prop:trivial} for the atom cover of the top, and list
its amalgamations. Then compute $a\Omega$ for $J_2$ and verify that the
result classifies $J_2$-closed sieves only.
\end{exercise}

\begin{exercise}\label{ch04:exr:twotopologies}
On $L_2$, let $J_a$ be generated by the single family $\{\{a\}\to\{a,b\}\}$
(and identities and the null cover), and $J_2$ as above. Determine which
representables are sheaves for $J_a$, whether $J_a\subseteq J_2$, and whether
the two modalities $M_{J_a}$ and $M_{J_2}$ are distinct, exhibiting a
separating presheaf if so. (This exercise previews the witness-presheaf
technique of the finite theory.)
\end{exercise}

\begin{exercise}\label{ch04:exr:leftexact}
Using the description $F^{+}(C)=\mathrm{Match}(S_{\min}(C),F)$ valid on
finite sites (Example~\ref{ch04:ex:plusworked}), verify by direct computation
that $(-)^{+}$ preserves finite products and equalizers on
$(L_2,J_2)$, and exhibit the natural isomorphism
$a(F\times G)\cong aF\times aG$ for $F=G=\Delta_X$.
\end{exercise}
